\section
{Glosario}
\label{cap:glosario}

El presente glosario reúne los términos técnicos y conceptuales de mayor 
recurrencia en esta investigación, organizados alfabéticamente. Cada 
definición incluye la fuente de referencia que le da sustento dentro del 
marco teórico del proyecto.

\begin{description}[leftmargin=0pt, labelwidth=\textwidth, 
                    labelsep=0pt, style=nextline]

% -------------------------------------------------------
\item[\textbf{Adherencia al tratamiento}]
Grado de continuidad con que un paciente realiza las actividades de una 
intervención a lo largo del tiempo. En el contexto de la estimulación 
cognitiva domiciliaria, la adherencia es un factor crítico para la 
efectividad del programa, ya que los beneficios dependen de la práctica 
sostenida \parencite{moret, calatayud2021}.

% -------------------------------------------------------
\item[\textbf{Alucinación (LLMs)}]
Fenómeno propio de los Grandes Modelos de Lenguaje por el cual el modelo 
genera respuestas lingüísticamente plausibles pero factualmente incorrectas 
o inventadas. Representa una limitación relevante en sistemas de salud 
digital donde la precisión de la información es crítica \parencite{moret}.

% -------------------------------------------------------
\item[\textbf{API (\textit{Application Programming Interface})}]
Interfaz de programación de aplicaciones: conjunto de protocolos y 
definiciones que permite a un programa de software comunicarse con otro. 
En este proyecto se utiliza la Gemini API de Google para integrar el 
módulo de PLN al motor de juego Unity \parencite{Ortega2024}.

% -------------------------------------------------------
\item[\textbf{Arquitectura por capas}]
Patrón de diseño de software que organiza un sistema en módulos 
independientes con responsabilidades bien delimitadas que se comunican 
mediante interfaces estandarizadas. Se identifican tres capas recurrentes 
en sistemas de salud digital: presentación e interacción, inteligencia 
artificial y datos \parencite{Goumopoulos2021, Karapapas2021}.

% -------------------------------------------------------
\item[\textbf{Asistente virtual (VA)}]
Sistema de software basado en procesamiento de lenguaje natural capaz de 
interactuar con el usuario mediante lenguaje natural hablado o escrito. 
En herramientas de estimulación cognitiva, el asistente virtual cumple 
funciones de instrucción, motivación y retroalimentación adaptativa 
\parencite{moret}.

% -------------------------------------------------------
\item[\textbf{Atención sostenida}]
Capacidad cognitiva para mantener el foco en una tarea o estímulo durante 
un período prolongado sin pérdida de rendimiento. Es una de las funciones 
cognitivas más afectadas en el Deterioro Cognitivo Leve y constituye una 
esfera central de estimulación en este proyecto \parencite{bermejo2021, 
AbdAlrazaq2022}.

% -------------------------------------------------------
\item[\textbf{Design Science Research (DSR)}]
Paradigma metodológico de investigación aplicada cuyo objetivo es diseñar, 
construir y evaluar artefactos tecnológicos que resuelvan problemas del 
mundo real. A diferencia de la investigación descriptiva, el DSR produce 
como resultado principal el artefacto mismo y la evidencia de su 
funcionamiento \parencite{Munoz2025}.

% -------------------------------------------------------
\item[\textbf{Deterioro Cognitivo Leve (DCL)}]
Etapa clínica intermedia entre el envejecimiento cognitivo normal y la 
demencia, caracterizada por un declive perceptible en una o más funciones 
cognitivas —memoria, atención o funciones ejecutivas— que no compromete 
significativamente la autonomía funcional del individuo. Representa una 
ventana de oportunidad para la intervención preventiva 
\parencite{Giovagnoli2024}.

% -------------------------------------------------------
\item[\textbf{Dispositivo háptico}]
Hardware que transmite sensaciones físicas al usuario —resistencia, 
vibración o guía de movimiento— como canal de retroalimentación en 
sistemas interactivos. Se clasifican en dispositivos anclados al suelo 
(\textit{grounded}), de mano (\textit{hand-held}) y vestibles 
(\textit{wearable}) \parencite{Adilkhanov2022}.

% -------------------------------------------------------
\item[\textbf{DSR}]
Véase \textit{Design Science Research}.

% -------------------------------------------------------
\item[\textbf{Estimulación cognitiva (EC)}]
Conjunto de actividades y estrategias diseñadas para activar y ejercitar 
sistemáticamente las funciones cognitivas del individuo, con propósito 
preventivo y de mantenimiento —a diferencia de la rehabilitación clínica, 
que busca recuperar funciones perdidas—. Su eficacia ha sido documentada 
en memoria de trabajo, funciones ejecutivas y reducción de indicadores 
de depresión \parencite{bermejo2021, calatayud2021, arteaga2024}.

% -------------------------------------------------------
\item[\textbf{Exergame}]
Videojuego que requiere actividad física del usuario como mecanismo 
principal de interacción, combinando ejercicio motor con tareas cognitivas. 
En el contexto del envejecimiento activo, los exergames son una modalidad 
de entrenamiento dual con evidencia empírica de efectividad 
\parencite{exergames2026}.

% -------------------------------------------------------
\item[\textbf{Firebase}]
Plataforma de desarrollo de aplicaciones en la nube de Google, utilizada 
en este proyecto como infraestructura para integrar la Gemini API al motor 
de juego Unity mediante comunicación segura con la capa de inteligencia 
artificial del sistema \parencite{Ortega2024}.

% -------------------------------------------------------
\item[\textbf{Flexibilidad cognitiva}]
Capacidad para adaptar el pensamiento y la conducta ante cambios en las 
reglas o condiciones de una tarea. Es un componente central de las 
funciones ejecutivas y una de las habilidades más afectadas en el DCL 
\parencite{AbdAlrazaq2022, Giovagnoli2024}.

% -------------------------------------------------------
\item[\textbf{Force Feedback (FFB)}]
Mecanismo de retroalimentación háptica que utiliza motores de torque para 
ejercer resistencia activa sobre el eje de un dispositivo, simulando 
fuerzas físicas del entorno virtual. Se distingue de la vibración ERM en 
que genera resistencia real al movimiento del usuario 
\parencite{Marchal2008}.

% -------------------------------------------------------
\item[\textbf{Funciones ejecutivas (FE)}]
Conjunto de procesos cognitivos de alto nivel que permiten planificar, 
organizar, iniciar y supervisar conductas orientadas a un objetivo; 
comprenden memoria de trabajo, flexibilidad cognitiva, inhibición de 
respuestas y toma de decisiones. Son particularmente vulnerables al 
deterioro en adultos mayores \parencite{AbdAlrazaq2022}.

% -------------------------------------------------------
\item[\textbf{Gamificación}]
Aplicación de mecánicas y elementos propios del diseño de videojuegos 
—puntuación, progresión de niveles, retroalimentación inmediata, 
recompensas— en contextos no lúdicos, con el objetivo de incrementar la 
motivación y el compromiso del usuario \parencite{abd, Ortega2024}.

% -------------------------------------------------------
\item[\textbf{HID (\textit{Human Interface Device})}]
Clase de dispositivos de entrada y salida de datos estándar que se 
comunican con un ordenador sin requerir controladores adicionales, 
utilizando el protocolo USB HID. Los volantes comerciales como el 
Kontorland AP5 operan bajo este protocolo \parencite{Adilkhanov2022}.

% -------------------------------------------------------
\item[\textbf{Háptica}]
Rama de la interacción humano-computadora que estudia el sentido del tacto 
y la percepción cinestésica como canal de comunicación entre el usuario y 
el sistema. Permite transmitir información mediante sensaciones físicas 
como resistencia, vibración o presión \parencite{Marchal2008, 
Adilkhanov2022}.

% -------------------------------------------------------
\item[\textbf{IHC (Interacción Humano-Computadora)}]
Disciplina que estudia el diseño, evaluación e implementación de sistemas 
de cómputo para uso humano, con énfasis en la experiencia del usuario. 
Sus principios guían el diseño de interfaces accesibles para adultos 
mayores \parencite{Peres2022SUS}.

% -------------------------------------------------------
\item[\textbf{KPI (\textit{Key Performance Indicator})}]
Indicador clave de rendimiento: métrica cuantificable definida para 
evaluar el grado de cumplimiento de un objetivo. En el presente prototipo, 
los KPIs miden eficacia (tareas completadas correctamente), eficiencia 
(tiempo por tarea) y satisfacción (puntuación SUS) \parencite{Peres2022SUS, 
Ortega2024}.

% -------------------------------------------------------
\item[\textbf{LLM (\textit{Large Language Model})}]
Gran Modelo de Lenguaje: modelo de inteligencia artificial entrenado sobre 
enormes corpus de texto mediante arquitecturas de transformadores con 
mecanismos de auto-atención, capaz de generar texto contextualmente 
coherente y mantener conversaciones fluidas en lenguaje natural. Entre los 
más utilizados en salud digital se encuentran GPT (OpenAI), Gemini 
(Google) y LLaMA (Meta) \parencite{moret, Hamdi2025}.

% -------------------------------------------------------
\item[\textbf{MCI (\textit{Mild Cognitive Impairment})}]
Término en inglés equivalente al Deterioro Cognitivo Leve (DCL). Véase 
\textit{Deterioro Cognitivo Leve}.

% -------------------------------------------------------
\item[\textbf{Memoria de trabajo}]
Sistema cognitivo de capacidad limitada que permite mantener y manipular 
información de forma temporal durante la realización de tareas complejas. 
Es uno de los dominios cognitivos más frecuentemente afectados en el DCL 
y más entrenados en programas de estimulación cognitiva 
\parencite{AbdAlrazaq2022, bermejo2021}.

% -------------------------------------------------------
\item[\textbf{MoCA (\textit{Montreal Cognitive Assessment})}]
Instrumento de tamizaje cognitivo de 30 puntos, diseñado específicamente 
para detectar el Deterioro Cognitivo Leve con mayor sensibilidad (87\%) 
que el MMSE. Evalúa memoria, atención, funciones ejecutivas, lenguaje, 
habilidades visuoespaciales y orientación. Existe versión validada en 
español \parencite{Giovagnoli2024}.

% -------------------------------------------------------
\item[\textbf{MMSE (\textit{Mini-Mental State Examination})}]
Instrumento de tamizaje cognitivo de 30 puntos ampliamente utilizado en 
la clínica. Presenta menor sensibilidad para DCL (72--94\%) que el MoCA 
y efecto techo en personas con alta escolaridad. Su acceso requiere 
licencia \parencite{Giovagnoli2024}.

% -------------------------------------------------------
\item[\textbf{Motor ERM (\textit{Eccentric Rotating Mass})}]
Motor de vibración de masa excéntrica rotatoria: actuador que genera 
vibración mediante el giro de una masa desequilibrada. Es el mecanismo 
háptico predominante en controladores comerciales de bajo costo como 
mandos de consola y volantes de entrada. Genera patrones de vibración 
programables sin ejercer resistencia activa sobre el eje del dispositivo 
\parencite{Adilkhanov2022, Gardinger2024}.

% -------------------------------------------------------
\item[\textbf{Motor de juego (\textit{game engine})}]
Plataforma de software que proporciona los componentes fundamentales para 
el desarrollo de aplicaciones interactivas: renderizado gráfico, gestión 
de físicas, manejo de entradas, audio y herramientas de scripting. En este 
proyecto se utiliza Unity como motor principal \parencite{Ortega2024}.

% -------------------------------------------------------
\item[\textbf{Orientación espacial}]
Capacidad cognitiva para ubicarse y desplazarse en el entorno, 
reconociendo la posición propia en relación con el espacio circundante. 
Es uno de los dominios cognitivos objeto de estimulación en este prototipo, 
trabajado mediante la navegación en laberintos \parencite{bermejo2021}.

% -------------------------------------------------------
\item[\textbf{PLN (Procesamiento de Lenguaje Natural)}]
Rama de la inteligencia artificial que estudia la interacción entre 
computadoras y el lenguaje humano, habilitando capacidades como 
comprensión, generación y traducción de texto o voz de forma automatizada. 
Es la base tecnológica del asistente conversacional de este prototipo 
\parencite{moret}.

% -------------------------------------------------------
\item[\textbf{Praxias visomotoras}]
Capacidad para ejecutar movimientos voluntarios aprendidos en respuesta 
a estímulos visuales, requiriendo coordinación entre la percepción visual 
y la respuesta motora. En este proyecto se activan mediante la coordinación 
simultánea de manos y pies durante la interacción con el volante y el pedal 
\parencite{bermejo2021, AbdAlrazaq2022}.

% -------------------------------------------------------
\item[\textbf{\textit{Prompt}}]
Instrucción o texto de entrada que se proporciona a un LLM para guiar su 
respuesta. En sistemas de salud digital, el uso de prompts estructurados 
y acotados es una estrategia para mitigar el fenómeno de alucinación y 
mantener la interacción dentro de contextos controlados y seguros 
\parencite{Ortega2024, moret}.

% -------------------------------------------------------
\item[\textbf{Privacidad por diseño (\textit{Privacy by Design})}]
Principio de arquitectura de software que integra la protección de datos 
personales desde las etapas iniciales del diseño del sistema, en lugar 
de añadirla como medida posterior. En sistemas de salud domiciliaria 
implica almacenar datos localmente y minimizar la transmisión de 
información identificable a servicios en la nube \parencite{moret, 
Munoz2025}.

% -------------------------------------------------------
\item[\textbf{Reserva cognitiva (RC)}]
Capacidad del cerebro para tolerar los efectos de la patología neurológica 
sin manifestar síntomas clínicos, actuando como mecanismo de compensación 
activa. Se forja a través de la educación, la actividad intelectual y el 
compromiso social a lo largo de la vida \parencite{Stern2022, 
Giovagnoli2024}.

% -------------------------------------------------------
\item[\textbf{Serious Game}]
Aplicación de software que utiliza mecánicas y elementos propios del 
diseño de videojuegos con un propósito primario distinto al entretenimiento: 
en el contexto de la salud, su objetivo es promover comportamientos 
terapéuticos, educativos o de rehabilitación, incorporando retroalimentación 
inmediata, progresión de dificultad y motivación intrínseca 
\parencite{abd, Ortega2024}.

% -------------------------------------------------------
\item[\textbf{SUS (\textit{System Usability Scale})}]
Escala de usabilidad percibida de diez ítems que genera una puntuación 
de 0 a 100. Una puntuación $\geq$ 68 indica usabilidad aceptable; 
puntuaciones $>$ 80.3 se clasifican como excelentes para aplicaciones de 
salud digital. Es el instrumento de evaluación de usabilidad seleccionado 
para la prueba piloto de este proyecto \parencite{Peres2022SUS}.

% -------------------------------------------------------
\item[\textbf{Tarea dual (\textit{dual-task})}]
Paradigma experimental y de entrenamiento cognitivo que exige la 
ejecución simultánea de dos tareas —una motriz y una cognitiva—, 
revelando déficits atencionales que permanecen ocultos bajo condiciones 
de tarea simple. El entrenamiento dual ha demostrado mejoras 
significativamente mayores en funciones ejecutivas que el entrenamiento 
cognitivo aislado \parencite{Voelcker2025DualTask, Li2023DualTask}.

% -------------------------------------------------------
\item[\textbf{Tareas ecológicamente válidas}]
Actividades de entrenamiento o evaluación cognitiva que simulan 
situaciones de la vida cotidiana reconocibles por el usuario. En el diseño 
de Serious Games para adultos mayores, la validez ecológica favorece la 
adherencia y la transferencia de las habilidades entrenadas al contexto 
real \parencite{Ortega2024, Gardinger2024}.

% -------------------------------------------------------
\item[\textbf{Unity}]
Motor de juego multiplataforma desarrollado por Unity Technologies, con 
soporte nativo para dispositivos HID, scripting en C\# y conexión con 
APIs externas. Es el motor de mayor adopción en proyectos de Serious Games 
de salud y rehabilitación cognitiva \parencite{Ortega2024, Ribeiro2024}.

% -------------------------------------------------------
\item[\textbf{Usabilidad}]
Grado en que un sistema permite a usuarios específicos alcanzar objetivos 
concretos con efectividad, eficiencia y satisfacción en un contexto de uso 
determinado. En herramientas para adultos mayores, la usabilidad es un 
criterio de diseño prioritario dada la menor experiencia de este grupo con 
interfaces digitales \parencite{Peres2022SUS, Ortega2024}.

% -------------------------------------------------------
\item[\textbf{Volante háptico de bajo costo}]
Controlador de simulación de conducción de precio accesible (aproximadamente 
\$20--\$40 USD) que integra motores ERM para retroalimentación vibrotáctil 
y se comunica mediante protocolos HID estándar (XInput/DInput). El modelo 
Kontorland AP5, utilizado en este prototipo, pertenece a esta categoría 
\parencite{Adilkhanov2022, Gardinger2024}.

% -------------------------------------------------------
\item[\textbf{XInput / DInput}]
Protocolos estándar de Windows para la comunicación entre controladores 
de juego y aplicaciones de software. XInput es el estándar moderno para 
controles Xbox-compatibles; DInput (DirectInput) es el protocolo heredado 
que soporta una gama más amplia de dispositivos HID, incluyendo volantes 
comerciales de bajo costo \parencite{Adilkhanov2022}.

\end{description}